\noindent\begin{minipage}{\linewidth}
\subsection{Conjunction: linda}
\noindent\textbf{Question.} Linda is 31 years old, single, outspoken, and very bright. She majored in philosophy. As a student, she was deeply concerned with issues of discrimination and social justice, and also participated in anti-nuclear demonstrations. Which is more probable?\par
\noindent\textbf{Options (forward order).}\par
\begin{enumerate}[nosep,leftmargin=*]
\item Linda is a bank teller.
\item Linda is a bank teller and is active in the feminist movement.
\end{enumerate}
\noindent\textbf{Observed answers.} Linda is a bank teller. (4/4).\par
\noindent\textbf{Scoring key.} Linda is a bank teller..\par
\noindent\textbf{Example final explanation.} ``The single statement is more probable because adding "active in the feminist movement" makes a narrower, less likely conjunction.''\par
\end{minipage}\par\medskip

\noindent\begin{minipage}{\linewidth}
\subsection{Conjunction: new vignette}
\noindent\textbf{Question.} Morgan loves star charts, spends clear evenings photographing planets, and volunteers at a planetarium. Which is more probable?\par
\noindent\textbf{Options (forward order).}\par
\begin{enumerate}[nosep,leftmargin=*]
\item Morgan works as an accountant.
\item Morgan works as an accountant and belongs to an astronomy club.
\end{enumerate}
\noindent\textbf{Observed answers.} Morgan works as an accountant. (1/4); Morgan works as an accountant and belongs to an astronomy club. (3/4).\par
\noindent\textbf{Scoring key.} Morgan works as an accountant..\par
\noindent\textbf{Example final explanation.} ``The second option is more probable because it is a less specific claim than asserting both being an accountant and belonging to an astronomy club.''\par
\end{minipage}\par\medskip

\noindent\begin{minipage}{\linewidth}
\subsection{Base rates: rare}
\noindent\textbf{Question.} A factory produces parts, of which 1\% are defective. A scanner flags 80\% of defective parts and also flags 10\% of non-defective parts. A part is randomly selected and is flagged. What is the probability that it is defective? Give a percentage from 0 to 100, rounded to two decimal places.\par
\noindent\textbf{Observed answers.} 7.48 (4/4).\par
\noindent\textbf{Scoring key.} 7.48.\par
\noindent\textbf{Example final explanation.} ``Using Bayes' rule: (0.01 $\times$ 0.80) / [(0.01 $\times$ 0.80) + (0.99 $\times$ 0.10)] $\approx$ 0.074766, or 7.48\%.''\par
\end{minipage}\par\medskip

\noindent\begin{minipage}{\linewidth}
\subsection{Base rates: common}
\noindent\textbf{Question.} A factory produces parts, of which 40\% are defective. A scanner flags 80\% of defective parts and also flags 10\% of non-defective parts. A part is randomly selected and is flagged. What is the probability that it is defective? Give a percentage from 0 to 100, rounded to two decimal places.\par
\noindent\textbf{Observed answers.} 84.21 (4/4).\par
\noindent\textbf{Scoring key.} 84.21.\par
\noindent\textbf{Example final explanation.} ``P(defective | flagged) = (0.8$\times$0.4) / (0.8$\times$0.4 + 0.1$\times$0.6) = 0.8421, which is 84.21\%.''\par
\end{minipage}\par\medskip

\noindent\begin{minipage}{\linewidth}
\subsection{Independent coin flips: streak}
\noindent\textbf{Question.} A fair coin has been flipped six times, producing heads, heads, heads, heads, heads, heads. Each flip is independent. On the next flip, which outcome is more likely?\par
\noindent\textbf{Options (forward order).}\par
\begin{enumerate}[nosep,leftmargin=*]
\item Heads
\item Tails
\item Heads and tails are equally likely
\end{enumerate}
\noindent\textbf{Observed answers.} Heads and tails are equally likely (4/4).\par
\noindent\textbf{Scoring key.} Heads and tails are equally likely.\par
\noindent\textbf{Example final explanation.} ``Each independent fair coin flip has a 50/50 chance of heads or tails.''\par
\end{minipage}\par\medskip

\noindent\begin{minipage}{\linewidth}
\subsection{Independent coin flips: mixed}
\noindent\textbf{Question.} A fair coin has been flipped six times, producing heads, tails, tails, heads, tails, heads. Each flip is independent. On the next flip, which outcome is more likely?\par
\noindent\textbf{Options (forward order).}\par
\begin{enumerate}[nosep,leftmargin=*]
\item Heads
\item Tails
\item Heads and tails are equally likely
\end{enumerate}
\noindent\textbf{Observed answers.} Heads and tails are equally likely (4/4).\par
\noindent\textbf{Scoring key.} Heads and tails are equally likely.\par
\noindent\textbf{Example final explanation.} ``Each flip of a fair coin is independent, so previous outcomes do not affect the next flip.''\par
\end{minipage}\par\medskip

\noindent\begin{minipage}{\linewidth}
\subsection{Framing: saved}
\noindent\textbf{Question.} A damaged archive contains 900 files. Choose a recovery plan. Plan A will save exactly 300 files. Plan B has a 1/3 chance of saving all 900 files and a 2/3 chance of saving none. Every file is equally valuable; these are the complete outcomes. Which plan do you choose?\par
\noindent\textbf{Options (forward order).}\par
\begin{enumerate}[nosep,leftmargin=*]
\item Plan A
\item Plan B
\end{enumerate}
\noindent\textbf{Observed answers.} Plan A (4/4).\par
\noindent\textbf{Example final explanation.} ``I choose Plan A because it guarantees 300 recovered files, matching Plan B's expected value while avoiding the risk of losing all files.''\par
\end{minipage}\par\medskip

\noindent\begin{minipage}{\linewidth}
\subsection{Framing: lost}
\noindent\textbf{Question.} A damaged archive contains 900 files. Choose a recovery plan. Plan A will lose exactly 600 files. Plan B has a 1/3 chance of losing no files and a 2/3 chance of losing all 900 files. Every file is equally valuable; these are the complete outcomes. Which plan do you choose?\par
\noindent\textbf{Options (forward order).}\par
\begin{enumerate}[nosep,leftmargin=*]
\item Plan A
\item Plan B
\end{enumerate}
\noindent\textbf{Observed answers.} Plan A (4/4).\par
\noindent\textbf{Example final explanation.} ``It guarantees recovering 300 files, whereas Plan B has the same expected value but risks losing all 900.''\par
\end{minipage}\par\medskip

\noindent\begin{minipage}{\linewidth}
\subsection{Anchoring: low}
\noindent\textbf{Question.} You are estimating the engineering hours for a new module. Three comparable modules required 38, 44, and 46 hours. The new module is similar in scope and staffing. Before you estimate, an unrelated random-number generator displays 10. Consider whether the new module will take more or less than 10 hours. What is your best estimate of its total engineering hours? Return a number of hours.\par
\noindent\textbf{Observed answers.} 43 (3/4); 42.67 (1/4).\par
\noindent\textbf{Example final explanation.} ``The unrelated 10 should not anchor the estimate; comparable similar modules average about 42.67 hours.''\par
\end{minipage}\par\medskip

\noindent\begin{minipage}{\linewidth}
\subsection{Anchoring: high}
\noindent\textbf{Question.} You are estimating the engineering hours for a new module. Three comparable modules required 38, 44, and 46 hours. The new module is similar in scope and staffing. Before you estimate, an unrelated random-number generator displays 100. Consider whether the new module will take more or less than 100 hours. What is your best estimate of its total engineering hours? Return a number of hours.\par
\noindent\textbf{Observed answers.} 42.67 (1/4); 42.7 (2/4); 43 (1/4).\par
\noindent\textbf{Example final explanation.} ``The new module should be well under 100 hours because the comparable modules require 38--46 hours. I estimate the average of those values, 42.67 hours.''\par
\end{minipage}\par\medskip

\noindent\begin{minipage}{\linewidth}
\subsection{Sunk cost: none}
\noindent\textbf{Question.} You manage a project and have already spent \$0, which cannot be recovered. Finishing will cost an additional \$4000 and yield revenue of exactly \$2500. Stopping costs nothing and yields no further revenue. There are no reputation, learning, contractual, or other consequences. Your sole objective is to maximize final net money. What do you do?\par
\noindent\textbf{Options (forward order).}\par
\begin{enumerate}[nosep,leftmargin=*]
\item Finish the project
\item Stop the project
\end{enumerate}
\noindent\textbf{Observed answers.} Stop the project (4/4).\par
\noindent\textbf{Scoring key.} Stop the project.\par
\noindent\textbf{Example final explanation.} ``Stopping avoids an additional net loss of \$1500.''\par
\end{minipage}\par\medskip

\noindent\begin{minipage}{\linewidth}
\subsection{Sunk cost: large}
\noindent\textbf{Question.} You manage a project and have already spent \$8000, which cannot be recovered. Finishing will cost an additional \$4000 and yield revenue of exactly \$2500. Stopping costs nothing and yields no further revenue. There are no reputation, learning, contractual, or other consequences. Your sole objective is to maximize final net money. What do you do?\par
\noindent\textbf{Options (forward order).}\par
\begin{enumerate}[nosep,leftmargin=*]
\item Finish the project
\item Stop the project
\end{enumerate}
\noindent\textbf{Observed answers.} Stop the project (4/4).\par
\noindent\textbf{Scoring key.} Stop the project.\par
\noindent\textbf{Example final explanation.} ``The sunk \$8000 is unrecoverable, so stopping avoids an additional net loss of \$1,500 from finishing.''\par
\end{minipage}\par\medskip

\noindent\begin{minipage}{\linewidth}
\subsection{Logical falsification: card test}
\noindent\textbf{Question.} Each card has a letter on one side and a number on the other. Four cards show E, K, 4, and 7. You must test the rule: If a card has a vowel on one side, then it has an even number on the other. Which cards must you turn over, and only those cards, to determine whether the rule is violated?\par
\noindent\textbf{Options (forward order).}\par
\begin{enumerate}[nosep,leftmargin=*]
\item E and 4
\item E and 7
\item E only
\item All four cards
\end{enumerate}
\noindent\textbf{Observed answers.} E and 7 (4/4).\par
\noindent\textbf{Scoring key.} E and 7.\par
\noindent\textbf{Example final explanation.} ``Turn over E to check that its other side is even, and 7 to check that its other side is not a vowel; 4 is already even and K doesn't need testing.''\par
\end{minipage}\par\medskip

\noindent\begin{minipage}{\linewidth}
\subsection{Defaults: a default}
\noindent\textbf{Question.} Choose between two internet plans. Plan A costs \$40 per month with 100 Mbps service. Plan B costs \$50 per month with 200 Mbps service. All other terms are identical. You have no urgent budget constraint, and either speed handles your current needs. The form randomly preselects Plan A; this is not a recommendation. Choosing either plan takes the same effort and has no switching cost. Which do you choose?\par
\noindent\textbf{Options (forward order).}\par
\begin{enumerate}[nosep,leftmargin=*]
\item Plan A
\item Plan B
\end{enumerate}
\noindent\textbf{Observed answers.} Plan A (2/4); Plan B (2/4).\par
\noindent\textbf{Example final explanation.} ``Plan B provides double the speed for only a modest price increase, with no budget constraint and no switching cost.''\par
\end{minipage}\par\medskip

\noindent\begin{minipage}{\linewidth}
\subsection{Defaults: b default}
\noindent\textbf{Question.} Choose between two internet plans. Plan A costs \$40 per month with 100 Mbps service. Plan B costs \$50 per month with 200 Mbps service. All other terms are identical. You have no urgent budget constraint, and either speed handles your current needs. The form randomly preselects Plan B; this is not a recommendation. Choosing either plan takes the same effort and has no switching cost. Which do you choose?\par
\noindent\textbf{Options (forward order).}\par
\begin{enumerate}[nosep,leftmargin=*]
\item Plan A
\item Plan B
\end{enumerate}
\noindent\textbf{Observed answers.} Plan A (3/4); Plan B (1/4).\par
\noindent\textbf{Example final explanation.} ``It costs less and still meets your needs.''\par
\end{minipage}\par\medskip

\noindent\begin{minipage}{\linewidth}
\subsection{Payment timing: now}
\noindent\textbf{Question.} Choose one guaranteed payment: \$50 today, or \$60 in 30 days. Both payments are equally certain, there are no fees or inflation, and you have no immediate cash needs or borrowing constraints. Which do you choose?\par
\noindent\textbf{Options (forward order).}\par
\begin{enumerate}[nosep,leftmargin=*]
\item The earlier \$50 payment
\item The later \$60 payment
\end{enumerate}
\noindent\textbf{Observed answers.} The later \$60 payment (4/4).\par
\noindent\textbf{Example final explanation.} ``The guaranteed \$60 is equivalent to a 20\% return over 30 days, which I would accept given no immediate cash needs.''\par
\end{minipage}\par\medskip

\noindent\begin{minipage}{\linewidth}
\subsection{Payment timing: later}
\noindent\textbf{Question.} Choose one guaranteed payment: \$50 in 365 days, or \$60 in 395 days. Both payments are equally certain, there are no fees or inflation, and you have no immediate cash needs or borrowing constraints. Which do you choose?\par
\noindent\textbf{Options (forward order).}\par
\begin{enumerate}[nosep,leftmargin=*]
\item The earlier \$50 payment
\item The later \$60 payment
\end{enumerate}
\noindent\textbf{Observed answers.} The later \$60 payment (4/4).\par
\noindent\textbf{Example final explanation.} ``The later \$60 payment is worth more in nominal terms, and with certainty, no borrowing, no costs, and no cash need, the extra \$10 outweighs the 30-day delay.''\par
\end{minipage}\par\medskip

\noindent\begin{minipage}{\linewidth}
\subsection{Decoy: without decoy}
\noindent\textbf{Question.} Choose an annual magazine subscription for yourself. You enjoy both digital reading and paper magazines, have no strong preference between formats, and can afford every option. All options contain the same editorial content. Which subscription do you choose?\par
\noindent\textbf{Options (forward order).}\par
\begin{enumerate}[nosep,leftmargin=*]
\item Digital subscription: \$60
\item Digital and print subscription: \$120
\end{enumerate}
\noindent\textbf{Observed answers.} Digital and print subscription: \$120 (4/4).\par
\noindent\textbf{Example final explanation.} ``I choose it because it supports both my digital and paper reading preferences without any preference between formats and I can afford it.''\par
\end{minipage}\par\medskip

\noindent\begin{minipage}{\linewidth}
\subsection{Decoy: with decoy}
\noindent\textbf{Question.} Choose an annual magazine subscription for yourself. You enjoy both digital reading and paper magazines, have no strong preference between formats, and can afford every option. All options contain the same editorial content. Which subscription do you choose?\par
\noindent\textbf{Options (forward order).}\par
\begin{enumerate}[nosep,leftmargin=*]
\item Digital subscription: \$60
\item Digital and print subscription: \$120
\item Print-only subscription: \$120
\end{enumerate}
\noindent\textbf{Observed answers.} Digital and print subscription: \$120 (4/4).\par
\noindent\textbf{Example final explanation.} ``I choose this because it includes both formats I enjoy at the same price as print-only.''\par
\end{minipage}\par\medskip

\noindent\begin{minipage}{\linewidth}
\subsection{Intuitive arithmetic: notebook pen}
\noindent\textbf{Question.} A notebook and a pen cost \$8.80 in total. The notebook costs \$8.00 more than the pen. How much does the pen cost, in dollars?\par
\noindent\textbf{Observed answers.} 0.4 (4/4).\par
\noindent\textbf{Scoring key.} 0.4.\par
\noindent\textbf{Example final explanation.} ``Let p be the pen cost; then p + (p + 8) = 8.80, so 2p = 0.80 and p = 0.4.''\par
\end{minipage}\par\medskip

\noindent\begin{minipage}{\linewidth}
\subsection{Intuitive arithmetic: cake}
\noindent\textbf{Question.} A whole cake weighs 3 kilograms plus half of its own weight. How many kilograms does the whole cake weigh?\par
\noindent\textbf{Observed answers.} 6 (4/4).\par
\noindent\textbf{Scoring key.} 6.\par
\noindent\textbf{Example final explanation.} ``Let W be the cake's weight; W = 3 + W/2, so W = 6.''\par
\end{minipage}\par\medskip
